% Source: https://artofproblemsolving.com/community/c3h3317333_girl_math_painting_frame
% Stable identity: aops-topic-3317333-q01
The frame width is $5$, so the long and short bases differ by $10$. Let their lengths be $L>S>0$, and let each outer side contain $2k+1$ bases. Then $L=S+10$ and $(k+1)L+kS=105$, giving $S=\dfrac{95-10k}{2k+1}$. Positivity requires $k\le9$. At $k=9$, the bases are $S=\dfrac{5}{19}$ and $L=\dfrac{195}{19}$; alternating them fits each outer side and the inner side of length 95, with the corner trapezoid legs meeting along the frame diagonals. Thus the bound is attainable, and the maximum count is $4(19)=\boxed{76}$.
\begin{asy}
import geometry;
usepackage("amsmath");
size(270);
real L=195/19.0,ss=5/19.0;pair O=(52.5,52.5);draw((0,0)--(105,0)--(105,105)--(0,105)--cycle);draw((5,5)--(100,5)--(100,100)--(5,100)--cycle);for(int k=0;k<4;++k){real x=0;for(int j=0;j<19;++j){real w=j%2==0?L:ss;real left=j%2==0?x+5:x-5;real right=j%2==0?x+w-5:x+w+5;draw(rotate(90*k,O)*((x,0)--(left,5)));draw(rotate(90*k,O)*((x+w,0)--(right,5)));x+=w;}}label("$105$",(52.5,0),S);label("$95$",(52.5,5),N);
\end{asy}
